Absolute B magnitude of NGC 2639:
\[
M_B = -9.95 \log v_{max} + 3.15 = -9.95 \log(324) + 3.15 = -21.83
\hfill(20\%)
\]
\[
\log R_{25} = -0.249 M_B - 4.00 = 1.4357
\;\;\Longrightarrow\;\;
R_{25} = 27.2678~\mathrm{pc} = 8.41 \times 10^{19}~\mathrm{cm}
\hfill(20\%)
\]
% sic: source says "pc" and "8.41e19 cm"; from the definition R25 is in kpc
% (27.2678 kpc = 8.41e22 cm), and the final mass below is consistent with kpc.
Mass estimation of NGC 2639:
\[
\mathcal{M} = \frac{v^2 R_{25}}{G}
= \frac{(324 \times 10^{5})^{2}\,(8.41 \times 10^{19})}{6.6726 \times 10^{-8}}
= 1.32 \times 10^{42}~\mathrm{kg}
= 6.62 \times 10^{11}~\mathcal{M}_\odot
\hfill(20\%)
\]
% sic: mixed CGS/SI units as in the original
The B magnitude of the Sun:
\[
M_{B\odot} = M_{V\odot} + (m_{B\odot} - m_{V\odot}) = 4.82 + 0.64 = 5.46
\hfill(20\%)
\]
meaning the Sun is fainter in the blue band relative to the A0 standard star.

Luminosity of NGC 2639 in solar luminosity ($L_\odot$):
\[
L = 10^{0.4(M_{B\odot} - M_B)} L_\odot
= 10^{0.4(5.46 + 21.83)} L_\odot
= 8.24 \times 10^{10}\,L_\odot
\hfill(20\%)
\]
